% Generated draft: requires mathematical and source review.
\documentclass[11pt]{article}
\usepackage[margin=25mm]{geometry}
\usepackage{fontspec}
\setmainfont{Latin Modern Roman}
\usepackage{amsmath,amssymb,mathtools,unicode-math}
\setmathfont{Latin Modern Math}
\usepackage{enumitem,needspace,xurl,hyperref}
\setlength{\parindent}{0pt}
\setlength{\parskip}{5pt}
\setlength{\emergencystretch}{3em}
\newcommand{\N}{\mathbb{N}}
\newcommand{\Z}{\mathbb{Z}}
\newcommand{\Q}{\mathbb{Q}}
\newcommand{\R}{\mathbb{R}}
\newcommand{\C}{\mathbb{C}}
\newcommand{\F}{\mathbb{F}}
\newcommand{\A}{\mathbb{A}}
\newcommand{\PP}{\mathbb{P}}
\newcommand{\E}{\mathbb{E}}
\newcommand{\eps}{\varepsilon}
\newcommand{\dd}{\,\mathrm{d}}
\begin{document}
\section*{UnsolvedMath ID: 30005930}
\textbf{Source problem code:} catalogue entry
\subsection{Definitions and mathematical statement}
A dimer configuration on a finite graph \(G=(V,E)\) is a perfect matching: a subset \(\tau\subseteq E\) such that every vertex belongs to exactly one edge of \(\tau\). Write \(\mathcal T(G)\) for the set of such configurations. Interpreting these configurations as geometric tilings requires a specified graph and geometric realization, which the supplied excerpt does not provide.

For a specified sequence of graphs \(G_k\) admitting perfect matchings, probability measures \(\mu_k\) on \(\mathcal T(G_k)\), and rescaled observables \(X_k:\mathcal T(G_k)\to\mathcal X\) taking values in a common metric space \((\mathcal X,d)\), a scaling-limit question asks whether the laws \((X_k)_*\mu_k\) converge weakly and, if so, how to characterize their limits. A deterministic limit shape \(x_*\in\mathcal X\) means that
\[
\forall\varepsilon>0,\qquad
\lim_{k\to\infty}\mu_k\!\left(\left\{\tau\in\mathcal T(G_k):d(X_k(\tau),x_*)>\varepsilon\right\}\right)=0.
\]
These definitions describe possible targets; they do not specify the missing model or assert that convergence holds.

\textbf{Source question, with unresolved scope.} Determine scaling limits and limit shapes for random dimer tilings, especially in dimensions greater than two, after specifying the regions, underlying graphs, probability laws, boundary conditions, rescaling, observables, and topology. The supplied cleaned statement asks about regions in \(\R^n\) as \(n\to\infty\), whereas the original excerpt writes \(R_n\) and discusses discrete vector fields, scaling limits, and a comparison of two and three dimensions. It is therefore unresolved whether the limiting parameter is the ambient dimension, a size or mesh parameter at fixed dimension, or both. No single precise asymptotic proposition is determined by this record.

The excerpt proposes a correspondence between tilings \(\tau\) and divergence-free discrete vector fields \(f_\tau\), with measurable divergence-free vector fields as scaling limits in a suitable topology. It does not supply the correspondence, field normalization, admissible limiting fields, topology, or boundary conditions. These must be specified before this proposed formulation can be made precise.

\textbf{Supplied comparison and qualifications.} The excerpt attributes a two-dimensional large-deviation principle to Cohn, Kenyon, and Propp, using uniform measures on tilings of \(R_n\), the supremum norm on corresponding height functions \(h\), and limiting boundary data expressed through height functions. It describes a rate function depending only on \(\nabla h\), represented by an integral of a local entropy function \(\operatorname{ent}_2\) with an explicit formula, and reports that analysis of this formula yields a unique limit shape for a given region and boundary conditions. The hypotheses and precise rate-function normalization are absent, so this description is retained as attributed background rather than an unrestricted theorem. The statement that the three-dimensional model is expected not to be solvable is an expectation in the supplied excerpt, not an established mathematical assertion.

The supplied dated assessment labels the problem partially solved: it reports a large-deviation theory and law of large numbers for fine-mesh three-dimensional dimers under controlled boundary data, while describing general tilings, fluctuations, and explicit entropy as remaining open targets. The controlled boundary hypotheses and the exact remaining questions are not supplied. This is a statement-only draft requiring human mathematical review.
\subsection{Short English statement}
The problem asks how random dimer tilings behave at large scales, particularly beyond two dimensions. The record does not resolve whether the parameter tending to infinity is dimension or spatial resolution, and it leaves the model, boundary data, and convergence topology unspecified. Its supplied partially solved label and claims about remaining questions are dated curation, not independently verified conclusions.
\subsection{Sources}
\begin{itemize}
\item Local registry: \texttt{data/unsolvedmath/problems.json}; permanent UnsolvedMath ID \(30005930\); exact source problem code \texttt{catalogue entry}. The supplied record belongs to the Oberwolfach Reports Open Problems collection and gives proposed year 2024; no individual proposer is identified.
\item Supplied report citation: \emph{Statistical Physics and Random Surfaces} (2024), Oberwolfach Reports, DOI 10.4171/owr/2024/25, pp.\ 1387--1444. \url{https://doi.org/10.4171/owr/2024/25}. The background describes these reports as documenting mathematical workshops and their open-problem discussions. No problem-specific page or named contribution is supplied.
\item Supplied literature source, labeled Chandgotia--Sheffield--Wolfram: \url{https://arxiv.org/abs/2304.08468}.
\item Supplied literature source, labeled Kenyon--Wolfram: \url{https://arxiv.org/abs/2506.12171}.
\item Supplied literature source, labeled OWR 2024/25: \url{https://doi.org/10.4171/owr/2024/25}.
\item The original excerpt attributes the two-dimensional large-deviation result to Cohn, Kenyon, and Propp, citing reference [2]. The bibliographic entry for [2] is not included in the supplied record.
\item Literature assessment date: \texttt{2026-08-22}. The partially solved status and accompanying assessment are supplied dated curation, described in the record as literature triage derived from OWR curation; they have not been independently verified here. The record explicitly requests independent verification of time-sensitive status and cited claims.
\item The registry labels its statement \texttt{corrected\_verified} and says that only formula delimiters, notation, grammar, and typography were repaired while preserving mathematical content and scope. This supplied label does not resolve the formulation ambiguities or constitute independent mathematical verification.
\end{itemize}
\end{document}
